% Options for packages loaded elsewhere
\PassOptionsToPackage{unicode}{hyperref}
\PassOptionsToPackage{hyphens}{url}
\documentclass[
  11pt,
]{article}
\usepackage{xcolor}
\usepackage[margin=1.1in]{geometry}
\usepackage{amsmath,amssymb}
\setcounter{secnumdepth}{-\maxdimen} % remove section numbering
\usepackage{iftex}
\ifPDFTeX
  \usepackage[T1]{fontenc}
  \usepackage[utf8]{inputenc}
  \usepackage{textcomp} % provide euro and other symbols
\else % if luatex or xetex
  \usepackage{unicode-math} % this also loads fontspec
  \defaultfontfeatures{Scale=MatchLowercase}
  \defaultfontfeatures[\rmfamily]{Ligatures=TeX,Scale=1}
\fi
\usepackage{lmodern}
\ifPDFTeX\else
  % xetex/luatex font selection
\fi
% Use upquote if available, for straight quotes in verbatim environments
\IfFileExists{upquote.sty}{\usepackage{upquote}}{}
\IfFileExists{microtype.sty}{% use microtype if available
  \usepackage[]{microtype}
  \UseMicrotypeSet[protrusion]{basicmath} % disable protrusion for tt fonts
}{}
\makeatletter
\@ifundefined{KOMAClassName}{% if non-KOMA class
  \IfFileExists{parskip.sty}{%
    \usepackage{parskip}
  }{% else
    \setlength{\parindent}{0pt}
    \setlength{\parskip}{6pt plus 2pt minus 1pt}}
}{% if KOMA class
  \KOMAoptions{parskip=half}}
\makeatother
\usepackage{longtable,booktabs,array}
\newcounter{none} % for unnumbered tables
\usepackage{calc} % for calculating minipage widths
% Correct order of tables after \paragraph or \subparagraph
\usepackage{etoolbox}
\makeatletter
\patchcmd\longtable{\par}{\if@noskipsec\mbox{}\fi\par}{}{}
\makeatother
% Allow footnotes in longtable head/foot
\IfFileExists{footnotehyper.sty}{\usepackage{footnotehyper}}{\usepackage{footnote}}
\makesavenoteenv{longtable}
\setlength{\emergencystretch}{3em} % prevent overfull lines
\providecommand{\tightlist}{%
  \setlength{\itemsep}{0pt}\setlength{\parskip}{0pt}}
\usepackage{amsmath}
\usepackage{amssymb}
\setcounter{secnumdepth}{-1}
\providecommand{\tightlist}{}
\usepackage{microtype}
\usepackage{bookmark}
\IfFileExists{xurl.sty}{\usepackage{xurl}}{} % add URL line breaks if available
\urlstyle{same}
\hypersetup{
  pdftitle={Collapse as the Present: George Ellis's Evolving Block Universe and Event Density's Committing Substrate},
  pdfauthor={Allen Proxmire},
  hidelinks,
  pdfcreator={LaTeX via pandoc}}

\title{Collapse as the Present: George Ellis's Evolving Block Universe
and Event Density's Committing Substrate}
\author{Allen Proxmire}
\date{June 2026}

\begin{document}
\maketitle

\subsection{Preamble --- what this article does and does not
claim}\label{preamble-what-this-article-does-and-does-not-claim}

\begin{enumerate}
\def\labelenumi{\arabic{enumi}.}
\tightlist
\item
  \textbf{This is a comparison and position paper, not a new
  derivation.} It maps an independent, established research program ---
  George Ellis's Evolving Block Universe (EBU) and its quantum
  refinement, the Crystallizing Block Universe (CBU) --- onto Event
  Density, and identifies where each supplies what the other lacks.
  Ellis's results are \emph{cited, not re-derived}; ED's are its
  standing conditional derivations, cited inline.
\item
  \textbf{Convergence, not endorsement.} Ellis has, to the author's
  knowledge, no acquaintance with Event Density. Every point of
  agreement below is \textbf{independent arrival at the same conclusion
  from a different starting point} (Ellis from general relativity and
  the measurement problem; ED from a discrete substrate), never Ellis's
  sanction of ED. The phrase used throughout is \emph{convergence}.
\item
  \textbf{ED reproduces the physics both frameworks share.} Nothing here
  revises quantum mechanics or general relativity. The ED-specific
  content invoked (the constructed substrate, the collapse scale
  \(M_{crit}\), the dissolution of the Past Hypothesis) carries its own
  standing tier, flagged per claim.
\item
  \textbf{The two things ED offers Ellis's program are tiered, not
  asserted as fact.} (a) A \emph{constructed} substrate whose commitment
  drives the block's growth is a conditional structural derivation from
  the thirteen ED primitives (\texttt{Paper\_087}), not a completed
  quantum-gravity theory. (b) A collapse mass-scale \(M_{crit}\) is an
  \emph{unconfirmed prediction} (\texttt{Paper\_060}), not a result.
  Neither closes Ellis's open problems; each is a candidate answer to
  one he names.
\item
  \textbf{Tiers marked per claim.} The core convergence
  (collapse-as-present \(\approx\) commitment) is \emph{interpretive and
  structural}; the substrate engine is a \emph{conditional derivation};
  \(M_{crit}\) is a \emph{prediction}; the Past-Hypothesis dissolution
  is a \emph{reframe} (\texttt{Paper\_HowTheoryCoarseGrainReality} §13);
  the preferred-frame comparison is a \emph{limited structural
  convergence} --- both reject Minkowski realism and carry preferred
  worldlines, but only ED is in the khronometric modified-gravity class,
  and the two disagree on the source of the frame (§6).
\end{enumerate}

\begin{center}\rule{0.5\linewidth}{0.5pt}\end{center}

\subsection{1. What Event Density is, and the class this runs
in}\label{what-event-density-is-and-the-class-this-runs-in}

Event Density is a discrete, relational substrate: a participation graph
whose channels carry a bandwidth \(b \ge 0\) and a phase, with one
process primitive --- \textbf{commitment} (P11), the irreversible act by
which an uncommitted participation becomes a definite, recorded fact,
writing the arrow of time into the law rather than onto it
(\texttt{Paper\_087}). Space and geometry are emergent: the metric is
the reading of the bandwidth field, \(g \sim 1/b\)
(\texttt{Paper\_MetricFromTheGraph}, \texttt{Paper\_GR-I}). Measurement
is not a separate postulate --- it \emph{is} commitment: the
basis-selecting, irreversible actualization of one outcome from a
superposed participation (Report §4). ED's gravitational class is
\textbf{khronometric}: two tensor modes plus one propagating scalar, the
\emph{khronon} \(u_\mu = \partial_\mu T/|\partial T|\), the arrow made
geometric (\texttt{Paper\_GR-II}), with luminal cones collapsing the
coupling space to one dimension \(\lambda_{local}\), on which
\(\alpha_{2} = 0\) exactly and \(\alpha_{1} = -4\lambda_{local}\),
suppressed \textasciitilde70 orders below current bounds
(\texttt{Paper\_GR-IV}).

Held in that frame, ED makes four claims that will matter here: (i) the
present is the live edge where the uncommitted becomes committed; (ii)
the future is genuinely open (uncommitted) and the past fixed
(recorded); (iii) the arrow is primitive, and the smooth time-reversible
laws are its coarse-grained shadow
(\texttt{Paper\_HowTheoryCoarseGrainReality}); (iv) collapse is real,
basis-selecting, and irreducible to unitary evolution.

Every one of those four is also a load-bearing claim of Ellis's EBU/CBU
--- reached from the opposite direction.

\subsection{2. Ellis's Evolving and Crystallizing Block
Universe}\label{elliss-evolving-and-crystallizing-block-universe}

Ellis (2006--2014; \emph{On the Flow of Time}; \emph{Time and Spacetime:
The Crystallizing Block Universe}, with Rothman; \emph{The Evolving
Block Universe and the Meshing Together of Times}) rejects both the
static block (eternalism, ``time is an illusion'') and presentism. His
ontology is a \textbf{growing} block: the past and present exist and are
fixed; the \textbf{future does not exist} --- it is a set of
possibilities, not a region. Spacetime is the standard four-dimensional
block \emph{minus the future}, with a moving growth-front. The present
is the growing edge, and it is ontologically special: it is where the
indefinite future becomes the definite past.

The mechanism is the paper's title-phrase: \textbf{collapse is the
present.} In Ellis's own italic hypothesis, the uncertainty of the
future changes to the certainty of the past when wave-function collapse
occurs --- ``all the time, everywhere.'' Collapse is not a measurement
\emph{problem}; it is time happening. This is an
\textbf{objective-collapse} position: Ellis treats reduction as a real,
time-irreversible physical process, extra to unitary quantum theory, and
explicitly rejects the three standard evasions --- the ensemble reading
(no ensemble without individual events), decoherence-as-sufficient (it
diagonalizes but does not select one outcome), and many-worlds (no
block, plus the measure and preferred-basis problems). The future is
open \textbf{ontologically}, from real quantum indeterminism amplified
to macro and cosmic scale (inflation, chaos), not merely epistemically.

The CBU is the quantum refinement: because of delayed-choice and
quantum-eraser evidence, the transition from indefinite to definite does
not occur on a clean simultaneity surface. The crystallization surface
is \emph{crinkly and lagging} --- some micro ``past'' events fix only
later --- so ``the past is fixed'' is strictly the macroscopic,
coarse-grained statement.

Two further commitments are essential for the comparison, and both are
exactly where ED has something to add:

\begin{itemize}
\tightlist
\item
  \textbf{The deep substrate is posited, not built.} Ellis holds that
  classical spacetime and its growth must emerge from an underlying
  quantum-gravity layer, that the growth is \emph{``presumably''}
  collapse of a quantum-gravity wave function, and that any correct
  quantum-gravity theory \emph{must} yield an EBU at the macro level ---
  a selection rule. He notes discrete programs (causal sets, spin foams)
  as compatible. He does not construct the layer.
\item
  \textbf{A theory of collapse is named as the missing piece.} Ellis
  states that the outstanding major issue for the program is a viable
  theory of \emph{when and where} collapse occurs; he does \textbf{not}
  adopt a Penrose--Diosi mass threshold, and treats the collapse scale
  as set contextually by the local environment, with no universal
  number.
\end{itemize}

\subsection{\texorpdfstring{3. The convergence: collapse-as-present
\(\approx\)
commitment}{3. The convergence: collapse-as-present  commitment}}\label{the-convergence-collapse-as-present-commitment}

The central result of this comparison is a near-identity of the two
frameworks' account of \emph{now}, reached independently.

{\def\LTcaptype{none} % do not increment counter
\begin{longtable}[]{@{}
  >{\raggedright\arraybackslash}p{(\linewidth - 4\tabcolsep) * \real{0.3333}}
  >{\raggedright\arraybackslash}p{(\linewidth - 4\tabcolsep) * \real{0.3333}}
  >{\raggedright\arraybackslash}p{(\linewidth - 4\tabcolsep) * \real{0.3333}}@{}}
\toprule\noalign{}
\begin{minipage}[b]{\linewidth}\raggedright
\end{minipage} & \begin{minipage}[b]{\linewidth}\raggedright
Ellis (EBU/CBU)
\end{minipage} & \begin{minipage}[b]{\linewidth}\raggedright
Event Density
\end{minipage} \\
\midrule\noalign{}
\endhead
\bottomrule\noalign{}
\endlastfoot
The present & the growing edge; indefinite \(\to\) definite & the live
edge; uncommitted \(\to\) committed (P11) \\
``Now'' mechanism & objective wave-function collapse \emph{is} the flow
of time & commitment \emph{is} the arrow (P11) \\
Future / past & future does not exist (open); past fixed/recorded &
future uncommitted (open); past committed/recorded \\
Arrow & fundamental, already in QM via measurement & primitive (P11);
reversibility is the coarse-grained shadow \\
Collapse & real, objective, extra to unitary QM; rejects MWI /
decoherence-alone / ensemble & commitment is a real primitive, not
derivable from unitary evolution; same three rejections \\
Spacetime & emergent from a discrete quantum-gravity layer whose
collapse drives growth & emergent from the discrete committing
substrate \\
\end{longtable}
}

Ellis's italic hypothesis --- collapse as the transition of an
indefinite future into a definite past, occurring everywhere continually
--- is, term for term, the coarse-grained description of P11. What Ellis
calls the \emph{crystallizing present}, ED calls the \emph{committing
edge}. The tier of this identification is \textbf{interpretive and
structural}: it is a mapping between two conditional readings of the
same quantum and relativistic physics, not a shared theorem. But its
significance is not diminished by that. Two derivations that begin as
far apart as general relativity and a discrete participation graph, and
that arrive at the same non-standard ontology --- real flow, open
future, fixed past, ontologically special present, objective collapse as
the mechanism of becoming --- constitute a \textbf{consilience} across
independent starting assumptions.

\subsection{4. What ED supplies that Ellis posits: the substrate
engine}\label{what-ed-supplies-that-ellis-posits-the-substrate-engine}

Ellis's EBU requires a Scale-0 layer whose collapse actualizes the
future at the present-boundary, and asserts that the correct
quantum-gravity theory must yield exactly that. Ellis describes the
\emph{shape} of this engine --- a discrete floor that commits, growing
spacetime at the present --- and states plainly that he does not have
it.

Event Density is a candidate for that engine, and it is
\emph{constructed} rather than posited. The committing substrate is the
Scale-0 layer; the emergent metric \(g \sim 1/b\) and the weak-field
Einstein equation are derived from it
(\texttt{Paper\_MetricFromTheGraph}, \texttt{Paper\_GR-I..IV}); the
growth of the block is the accumulation of commitments (P11); the
``actualization of the future at the present-boundary'' is
measurement-as-commitment (Report §4). ED even meets Ellis's stated
compatibility class --- a discrete, causal, becoming-based substrate in
the lineage he names (causal sets, spin foams).

The honest tier: ED is a \textbf{conditional structural derivation from
thirteen posited primitives}, not a completed quantum gravity, and the
full nonlinear Einstein equation is not derived from the primitives
(\texttt{Paper\_GR-I} delivers the weak field and the khronometric
class). \emph{(A separate thermodynamic route,
\texttt{Paper\_KhronometricEquationOfState\_Jacobson}, recovers the
Einstein equation in its equation-of-state} form \emph{with the aether
kept among the sources; that is a literature-grounded result, not a
from-primitives derivation of the full nonlinear field equations, which
remains open.)} So ED does not \emph{close} Ellis's Scale-0 problem. It
offers a specific, constructed candidate for the layer his program
requires --- which is more than the selection rule (``QG must yield an
EBU'') currently provides.

\subsection{\texorpdfstring{5. The collapse scale Ellis lacks:
\(M_{crit}\)}{5. The collapse scale Ellis lacks: M\_\{crit\}}}\label{the-collapse-scale-ellis-lacks-m_crit}

Ellis names a viable theory of \emph{when and where} collapse occurs as
an outstanding issue --- and, importantly, declines a universal collapse
scale, treating the trigger as set contextually by the local
environment. Event Density goes further, and commits to a scale. If
commitment carries a critical multiplicity, there is a
quantum-to-classical crossover \(M_{crit}\) (\texttt{Paper\_060}): a
mass above which multi-channel participation can no longer be maintained
and the substrate must commit. The associated matter-wave band,
\textasciitilde140--250 kDa, is \textbf{inherited from empirical
extrapolation} (\texttt{Paper\_056}), not ED-derived; what ED
contributes is the \emph{mechanism} --- that a sharp, substrate-forced
individuation boundary exists in that regime --- and interferometry is
currently approaching it.

The tier is explicit: \(M_{crit}\) is an \textbf{unconfirmed
prediction}, not a result. It is a candidate answer to the ``when does
collapse happen'' question Ellis leaves open --- and a \emph{bolder} one
than he endorsed: where Ellis treats the trigger as contextual with no
universal number, ED commits to a universal scale, and is falsifiable
for exactly that reason. Where Ellis's CBU is \emph{designed to
accommodate} the existing collapse and delayed-choice experiments
(Guerlin et al.~2007; the Kim et al.~quantum eraser; Wheeler
delayed-choice), ED adds a falsifiable number: a boundary that a
static-block or purely-unitary account says should not exist, and that
objective collapse says should. This is the one place the two programs,
taken together, become sharper than either alone --- a shared ontology
plus a testable edge that goes beyond Ellis's own stance.

A note on the CBU's crinkly present. Ellis's delayed-choice lag ---
micro ``past'' events crystallizing later --- has a natural ED reading:
commitment occurs at the actual interaction that fixes a degree of
freedom, which can be later than the naïve past-surface. ED reproduces
the delayed-choice results (it reproduces the Born rule and
no-signaling, A1 capacity-zero), and the interpretive picture --- fixing
happens at the commit, not before --- is the same structure as Ellis's
lagging crystallization surface.

\subsection{6. Preferred frames: a narrow convergence and a real
divergence}\label{preferred-frames-a-narrow-convergence-and-a-real-divergence}

The hardest problem for any \emph{becoming} ontology is relativity's
lack of a preferred simultaneity. Both frameworks confront it, and both
end up carrying preferred worldlines while treating Lorentz invariance
as effective --- but they do so in \emph{different theories}, and, as it
turns out, they disagree on the direction of the dependence.

Ellis rejects physical realism about Minkowski spacetime (no real vacuum
exists; the cosmic background radiation permeates everywhere), so the
\emph{real} universe carries geometrically preferred worldlines
everywhere: the timelike \textbf{Ricci eigenlines}, the local average
rest frame of matter. Crucially, these are a foliation \emph{chosen
within} the standard Einstein field equations --- an ADM gauge / broken
solution-symmetry (like the CMB rest frame), \textbf{not} a modification
of gravity. Ellis posits no dynamical aether field and no
Lorentz-violating couplings. The growth is \emph{pointwise along
worldlines} by proper time, needing no global simultaneity.

Event Density also carries a preferred structure and also needs no
global ``now'' (commitment is local and pointwise, exactly Ellis's
pointwise-along-worldlines growth). But ED locates it in the
\textbf{commitment order} --- the khronon
\(u_\mu = \partial_\mu T/|\partial T|\), the arrow made geometric
(\texttt{Paper\_GR-II}) --- and this is a \emph{genuine} khronometric
(Einstein-aether-class) \textbf{modification} of gravity, with computed
preferred-frame parameters (\(\alpha_{2} = 0\) exactly,
\(\alpha_{1} = -4\lambda_{local}\), \textasciitilde70 orders suppressed;
\texttt{Paper\_GR-IV}). ED, not Ellis, lives in the khronometric family.

So the comparison must be stated precisely, because it cuts both ways:

\begin{itemize}
\tightlist
\item
  \textbf{The convergence is real but narrow.} Both reject Minkowski
  realism, both carry geometrically preferred worldlines in the real
  matter-filled universe, and both keep Lorentz invariance as an
  \emph{effective} symmetry over a deeper order. That is a genuine
  structural agreement about what taking becoming seriously requires of
  the relativistic structure.
\item
  \textbf{The theories are not the same.} Ellis stays inside unmodified
  GR (a preferred foliation); ED is a Lorentz-violating modified gravity
  (a dynamical khronon). Only ED is in the khronometric /
  Einstein-aether class
  (\texttt{Paper\_KhronometricEquationOfState\_Jacobson},
  \texttt{Paper\_GR-IV}).
\item
  \textbf{They disagree on the direction of the frame.} ED derives the
  preferred frame \emph{from} the commitment order (arrow \(\to\)
  frame). Ellis explicitly argues the reverse --- that nothing inherent
  in collapse supplies a frame, and that the local matter context
  defines the preferred frame, which then sets the context for collapse
  (matter \(\to\) frame \(\to\) collapse; \emph{Meshing} §3.2). This is
  a \textbf{substantive disagreement}, not a shared mechanism.
\end{itemize}

The honest reading, then: a shared \emph{instinct} --- that becoming
forces preferred worldlines and an effective, emergent Lorentz symmetry
--- not a shared class and not a shared mechanism. ED goes further than
Ellis, into a testable modified-gravity khronon, and in a direction
Ellis argued against.

\subsection{7. Where ED and Ellis diverge ---
honestly}\label{where-ed-and-ellis-diverge-honestly}

The convergence is not total, and the differences are real:

\begin{itemize}
\tightlist
\item
  \textbf{The Past Hypothesis.} To align the \emph{thermodynamic} arrow
  with the fundamental direction of time, Ellis retains a fine-tuned
  low-entropy initial condition (a ``past condition''). ED does
  \textbf{not} require it: with irreversibility primitive, the
  reversible laws are the lossy average and the second law is
  bookkeeping, so no special initial condition is imported to explain
  the arrow (\texttt{Paper\_HowTheoryCoarseGrainReality} §13). Tier:
  this ED move is a \textbf{reframe} --- it identifies low initial
  entropy with the fewest early commitments (itself a substantive
  assumption), not a computed derivation of the initial-entropy value
  --- but it is a genuine economy relative to Ellis's account.
\item
  \textbf{The preferred-frame locus} (§6): matter (Ellis) versus the
  commitment order (ED). Empirically both live in the khronometric
  class; the difference is which physical structure supplies the
  foliation.
\item
  \textbf{The integrability of the meshing.} Ellis asserts, on
  observational rather than derivational grounds, that the local proper
  times mesh into a single growing block (and flags the
  timelike-crystallization-surface case as unfinished). ED grounds the
  corresponding fact differently --- the emergent metric and its
  locality come from the graph --- but neither framework has a closed
  proof that the meshing is forced; it is a shared open seam.
\end{itemize}

\subsection{8. Honest tiers}\label{honest-tiers}

\begin{itemize}
\tightlist
\item
  \textbf{Collapse-as-present \(\approx\) commitment} (§3):
  \emph{interpretive / structural} --- a convergence of two independent
  readings, not a shared theorem. Load-bearing as consilience, not as
  proof.
\item
  \textbf{ED as the substrate engine} (§4): \emph{conditional structural
  derivation} from the thirteen primitives; the weak-field metric and
  khronometric class are delivered, the full nonlinear Einstein equation
  is not. A constructed candidate for Ellis's posited Scale-0, not a
  completed quantum gravity.
\item
  \textbf{\(M_{crit}\)} (§5): \emph{unconfirmed prediction}; the ED
  answer to the ``when does collapse happen'' question Ellis names as
  open.
\item
  \textbf{Past-Hypothesis dissolution} (§7): \emph{reframe}; an honest
  economy, not a computed initial-entropy value.
\item
  \textbf{Preferred-frame comparison} (§6): \emph{limited structural
  convergence} --- both reject Minkowski realism and carry preferred
  worldlines, and both treat Lorentz invariance as effective. But only
  ED is a khronometric modified-gravity theory, and Ellis explicitly
  locates the frame in matter and argues \emph{against} ED's
  arrow-as-frame direction. A shared instinct, not a shared class or a
  shared mechanism.
\item
  \textbf{Convergence, not endorsement}: throughout. Ellis has no
  knowledge of ED; the agreement is independent arrival.
\end{itemize}

\subsection{9. Position statement}\label{position-statement}

Event Density is not a rival to the Evolving Block Universe; it is a
candidate \textbf{engine} for it. Ellis, from general relativity and the
measurement problem, arrives at a growing block whose present is
objective wave-function collapse, whose future does not exist, and whose
deep basis must be a discrete becoming-substrate he describes but does
not build. Among the open problems he names, two are ones Event Density
speaks to directly: that substrate, and a theory of \emph{when} collapse
occurs. (Others he lists --- the timelike-surface initial-value problem,
and the meshing of local proper times into one block --- ED does not
resolve; see §7.) Event Density, from a discrete committing substrate,
arrives at the same ontology --- present as the committing edge, open
future, fixed past, objective collapse as the arrow --- and offers, at
their stated tiers, a \emph{constructed} substrate and a
\emph{predicted} collapse scale.

The honest claim is therefore narrow and, precisely because it is
narrow, strong: two frameworks built from opposite ends of physics
converge on a non-standard account of time that most of the field
rejects, and the discrete framework supplies two of the pieces the
relativistic one names as missing. That is convergence plus candidate
completion, not proof --- and it is the appropriate relationship between
an outsider's substrate program and one of the field's most serious
defenses of the reality of time.

\begin{center}\rule{0.5\linewidth}{0.5pt}\end{center}

\subsection{Cross-references}\label{cross-references}

\begin{itemize}
\tightlist
\item
  \textbf{Ellis's program:} \emph{On the Flow of Time} (Ellis, FQXi
  2008); \emph{Time and Spacetime: The Crystallizing Block Universe}
  (Ellis \& Rothman, arXiv:0912.0808); \emph{The Evolving Block Universe
  and the Meshing Together of Times} (Ellis).
\item
  \textbf{Substrate \& commitment:}
  \texttt{physics-papers/foundations/Paper\_087\_13Primitives.md} (P11);
  Report §4 (measurement as commitment).
\item
  \textbf{Emergent metric \& gravitational class:}
  \texttt{Paper\_MetricFromTheGraph\_ForcedTo3D.md};
  \texttt{Paper\_GR-I..IV}; the khronon \texttt{Paper\_GR-II}; the
  khronometric couplings \texttt{Paper\_GR-IV}; the companion
  \texttt{Paper\_KhronometricEquationOfState\_Jacobson.md}.
\item
  \textbf{The arrow as coarse-graining, and the Past-Hypothesis
  dissolution:}
  \texttt{physics-papers/substrate-evaluation/Paper\_HowTheoryCoarseGrainReality.md}
  (§13); companion \texttt{Paper\_ConservationColumn.md}.
\item
  \textbf{The collapse scale:}
  \texttt{physics-papers/q-compute/Paper\_060\_Mcrit\_Unification.md};
  the matter-wave prediction, \texttt{physics-papers/predictions/}.
\item
  \textbf{No-signalling / commitment-locality:} the A1 capacity-zero
  result.
\end{itemize}

\end{document}
